% Source: 11 - Tue Oct 20 2026.pdf, page 6. Content remains in source order.
\sourcepage{11}{6}
is a poly-time decision algorithm for HAMILTON, which contradicts $\classP\ne\classNP$.

\subsection{The general gap technique}
The proof implements the general GAP technique to prove inapproximability results:

\textbf{GIVEN:} minimization problem $\Pi_m$ for which I want to prove $\rho(n)$-inapproximability:

\textbf{DETERMINE:}
\begin{itemize}
\item decision problem $\Pi_d\in\classNPC$;
\item ptc function $f(x)$ transforming instances of $\Pi_d$ into instances of $\Pi_m$ such that:
\item Let $c:\mathcal S\mapsto\mathbb N^+$ be the cost function of $\Pi_m$: there exist a function $K(n)$ such that
\begin{enumerate}
\item $x\in L_{\Pi_d}\Rightarrow c(s^*(f(x)))\leq K(|f(x)|)$;
\item $x\notin L_{\Pi_d}\Rightarrow c(s^*(f(x)))>\rho(|f(x)|)\cdot K(|f(x)|)$.
\end{enumerate}
\item Then if $\classP\ne\classNP$ there cannot be a poly-time $\rho(n)$-approximation algorithm for $\Pi_m$.
\end{itemize}
\textbf{PROOF} If $A_{\Pi_m}^\rho$ existed:
\begin{notesalgorithm}
\textbf{DECIDE ($x$)}
\begin{notespseudocode}
\pseudocodeline{0}{$s\gets A_{\Pi_m}^\rho(A_f(x))$}
\pseudocodeline{0}{\textbf{if} $c(s)\leq\rho(|f(x)|)\cdot K(f(x))$}
\pseudocodeline{1}{\textbf{then return} $1$}
\pseudocodeline{1}{\textbf{else return} $0$}
\end{notespseudocode}
\end{notesalgorithm}
would decide $\Pi_d$ in polynomial time.
